% =====================================================================
% FONTES DEPENDENTES EM MALHAS --- NIVEL 4
% Questoes integradas ao bloco de Desafio do Capitulo 9.
% =====================================================================

\begin{questao}[4]{}
No circuito de três malhas, a fonte dependente da terceira malha vale
$2{,}5I_1$ volts e impulsiona $I_3$. Determine $I_1$, $I_2$ e $I_3$.

\circuitoq{%
\begin{circuitikz}[american,scale=.84,transform shape,line width=.8pt]
\draw (0,0) to[\fonteV,l^={$\SI{20}{\volt}$}] (0,3.2)
      to[R,l={$\SI{5}{\ohm}$}] (2.8,3.2) coordinate(t1);
\draw (t1) to[R,l_={$\SI{5}{\ohm}$}] (2.8,0) coordinate(b1) -- (0,0);
\draw (t1) to[R,l={$\SI{5}{\ohm}$}] (5.6,3.2) coordinate(t2);
\draw (t2) to[R,l_={$\SI{5}{\ohm}$}] (5.6,0) coordinate(b2) -- (b1);
\draw (t2) to[R,l={$\SI{5}{\ohm}$}] (8.4,3.2)
      to[cvsource,l={$2{,}5I_1$}] (8.4,0) -- (b2);
\foreach \x/\n in {1.25/1,4.15/2,7.05/3}{%
  \draw[->,loopcol,line width=1.05pt] (\x,1.00) arc (0:310:.40);
  \node[loopcol,fill=white,inner sep=1pt] at (\x,1.00){\scriptsize$I_{\n}$};}
\end{circuitikz}}

\resposta{$I_1=\dfrac{8}{3}\,\si{\ampere}$; $I_2=I_3=\dfrac{4}{3}\,\si{\ampere}$}
\resolucao{As três equações são
\[
10I_1-5I_2=20,
\]
\[
-5I_1+15I_2-5I_3=0,
\]
\[
-5I_2+10I_3=2{,}5I_1.
\]
A solução do sistema fornece $I_1=8/3$ A e $I_2=I_3=4/3$ A.}
\end{questao}

\begin{questao}[4]{}
A fonte de corrente dependente está entre as malhas 2 e 3 e impõe
$I_3-I_2=0{,}5I_1$. Determine as três correntes usando uma supermalha envolvendo
as malhas 2 e 3.

\circuitoq{%
\begin{circuitikz}[american,scale=.84,transform shape,line width=.8pt]
\draw (0,0) to[\fonteV,l^={$\SI{30}{\volt}$}] (0,3.2)
      to[R,l={$\SI{5}{\ohm}$}] (2.8,3.2) coordinate(t1);
\draw (t1) to[R,l_={$\SI{5}{\ohm}$}] (2.8,0) coordinate(b1) -- (0,0);
\draw (t1) to[R,l={$\SI{5}{\ohm}$}] (5.6,3.2) coordinate(t2);
\draw (t2) to[cisource,l={$0{,}5I_1$},invert] (5.6,0) coordinate(b2) -- (b1);
\draw (t2) to[R,l={$\SI{10}{\ohm}$}] (8.4,3.2) -- (8.4,0) -- (b2);
\foreach \x/\n in {1.25/1,4.15/2,7.05/3}{%
  \draw[->,loopcol,line width=1.05pt] (\x,1.00) arc (0:310:.40);
  \node[loopcol,fill=white,inner sep=1pt] at (\x,1.00){\scriptsize$I_{\n}$};}
\draw[dashed,laranja,line width=.9pt,rounded corners=7pt]
      (2.25,-.55) rectangle (8.95,4.05);
\node[laranja,font=\footnotesize\sffamily,fill=white,inner sep=1.5pt]
      at (5.6,-.85) {supermalha 2--3};
\end{circuitikz}}

\resposta{$I_1=\SI{3}{\ampere}$; $I_2=0$; $I_3=\SI{1.5}{\ampere}$}
\resolucao{A malha 1 fornece $10I_1-5I_2=30$. A supermalha 2--3 fornece
\[
-5I_1+10I_2+10I_3=0,
\]
e a fonte impõe $I_3-I_2=0{,}5I_1$. Resolvendo o sistema,
$I_1=\SI{3}{\ampere}$, $I_2=0$ e $I_3=\SI{1.5}{\ampere}$.}
\end{questao}

\begin{questao}[4]{}
A fonte dependente vale $kI_1$ volts. Determine $k$ para que
$I_2=\SI{2}{\ampere}$ e calcule $I_1$.

\circuitoq{%
\begin{circuitikz}[american,scale=.92,transform shape,line width=.8pt]
\draw (0,0) to[\fonteV,l^={$\SI{20}{\volt}$}] (0,3.2)
      to[R,l={$\SI{5}{\ohm}$}] (3.2,3.2) coordinate(t);
\draw (t) to[R,l={$\SI{5}{\ohm}$}] (3.2,0) coordinate(b) -- (0,0);
\draw (t) to[R,l={$\SI{5}{\ohm}$}] (6.4,3.2)
      to[cvsource,l={$kI_1$}] (6.4,0) -- (b);
\draw[->,loopcol,line width=1.2pt] (1.75,1.00) arc (0:310:.46);
\node[loopcol,fill=white,inner sep=1pt] at (1.75,1.00){\footnotesize$I_1$};
\draw[->,loopcol,line width=1.2pt] (4.75,1.00) arc (0:310:.46);
\node[loopcol,fill=white,inner sep=1pt] at (4.75,1.00){\footnotesize$I_2$};
\end{circuitikz}}

\resposta{$I_1=\SI{3}{\ampere}$; $k=\dfrac{5}{3}\,\ohm\approx\SI{1.67}{\ohm}$}
\resolucao{A malha 1 dá $10I_1-5I_2=20$. Impondo $I_2=2$ A,
$I_1=3$ A. Na malha 2,
\[
-5I_1+10I_2=kI_1.
\]
Logo $-15+20=3k$, e portanto $k=5/3\,\ohm$.}
\end{questao}

\begin{questao}[4]{}
A fonte de tensão dependente está no ramo comum e possui tensão, de cima para
baixo, $v_d=2I_1$. Determine $I_1$, $I_2$, a corrente real na fonte dependente e
sua potência, indicando se ela absorve ou fornece energia.

\circuitoq{%
\begin{circuitikz}[american,scale=.92,transform shape,line width=.8pt]
\draw (0,0) to[\fonteV,l^={$\SI{20}{\volt}$}] (0,3.2)
      to[R,l={$\SI{4}{\ohm}$}] (3.2,3.2) coordinate(t);
\draw (t) to[cvsource,l_={$v_d$}] (3.2,0) coordinate(b) -- (0,0);
\draw (t) to[R,l={$\SI{4}{\ohm}$}] (6.4,3.2) -- (6.4,0) -- (b);
\draw[->,loopcol,line width=1.2pt] (1.75,1.00) arc (0:310:.46);
\node[loopcol,fill=white,inner sep=1pt] at (1.75,1.00){\footnotesize$I_1$};
\draw[->,loopcol,line width=1.2pt] (4.75,1.00) arc (0:310:.46);
\node[loopcol,fill=white,inner sep=1pt] at (4.75,1.00){\footnotesize$I_2$};
\end{circuitikz}}

\resposta{$I_1=\dfrac{10}{3}\,\si{\ampere}$; $I_2=\dfrac{5}{3}\,\si{\ampere}$;
$i_d=\dfrac{5}{3}\,\si{\ampere}$ para baixo; $P_d=\dfrac{100}{9}\,\si{\watt}$ absorvidos}
\resolucao{Na malha 1, $4I_1+v_d=20$ e $v_d=2I_1$, portanto
$I_1=10/3$ A. Na malha 2, $4I_2-v_d=0$, então $I_2=5/3$ A. A corrente real no
ramo comum é $i_d=I_1-I_2=5/3$ A para baixo. Como $v_d=20/3$ V e a corrente
entra no terminal positivo superior da fonte dependente,
\[
P_d=v_di_d=\frac{20}{3}\frac{5}{3}=\frac{100}{9}\,\si{\watt},
\]
logo a fonte dependente \textbf{absorve} potência.}
\end{questao}
